\section{Discussion}
The combined results establish a reliability-aware active-learning workflow in which numerical trust remains part of each physical observation and exact BdG calculations retain label authority. The two-dimensional thermodynamic and configured topology audits establish the low-dimensional results, while the three-dimensional extension demonstrates why selective exact sampling becomes more consequential as the domain grows. That extension is deliberately bounded by its finite sampled interval, discrete exact support, and four fixed-$\mu$ slices; these data reveal structure but do not define a continuous topology surface or complete atlas.

The matched Stage-V audit likewise gives no overall acquisition winner: V3 has slightly denser topology-surface support ($p_{95}=0.01893$ versus $0.02103$), whereas V4 has lower sequential topology surprise ($0.0925$ versus $0.1189$). V4 does not lower the common offline graph-uncertainty functional ($0.01217$ versus $0.01114$), and none of its four candidate nodes has validated local same-identity support. Learned acquisition is therefore an operational exploratory mechanism, not evidence of superiority or autonomous new-phase discovery.

As a future direction, the Stage-VI studies indicate that thermodynamic state and $|q|$ are learnable and that continuous Pfaffian quantities become learnable after conditioning on trusted gapped superconductivity. The frozen export, however, lacks fixed-$k$ spectra, a density-of-states proxy, and validated low-energy spectral moments, and its replay uncertainty is not uniformly monotonic. A ground-up physical world model must therefore validate those descriptors and calibrated uncertainty before production expected-information-gain acquisition can be authorized.
